\section{Proof of abundance after nonvanishing}\label{sec:completion}

\begin{proof}[Proof of Theorem~\ref{thm:main}]
If \(\kappa(X,D)\geq1\), Proposition~\ref{prop:positive-kappa} gives
semiampleness on the original \(X\), by descent of the generated actual
Cartier multiple.  It remains to treat \(\kappa(X,D)=0\).

Choose a positive integer \(m_0\) for which the actual line
\(\mathscr L=\OO_X(m_0D)\) is invertible and has a nonzero section
\(s_0\).  A connected normal space is irreducible, so this section is
locally a nonzerodivisor.  Let
\[
 M=\frac1{m_0}\operatorname{div}_{\mathscr L}(s_0).
\]
It is the normalized effective rational \(\Q\)-Cartier divisor of the
plurisection, and the chosen line and section give the actual equivalence
\(M\sim_\Q D\).

Suppose \(M\ne0\).  Proposition~\ref{prop:supported-model} produces a
normal ordinary \(\Q\)-factorial compact K\"ahler dlt fourfold \((V,B)\)
with effective rational boundary, analytically nef actual adjoint
\(A=K_V+B\), and a nonzero effective rational \(\Q\)-Cartier divisor
\(P\) such that
\[
 A\sim_\Q P,\qquad \Supp P=\Supp\floor B,\qquad \kappa(V,A)=0.
\]
The space \(V\) is irreducible, being bimeromorphic to \(X\), and has
the special projective resolution in Lemma~\ref{lem:special-resolution}.
Theorem~\ref{thm:floor-generation} therefore makes the actual restricted
adjoint semiample on the entire reduced floor \(S=\floor B\).  This
floor is nonempty because \(P\ne0\).  Every hypothesis of
Theorem~\ref{thm:supported-lifting} is now satisfied, so it gives
\(\kappa(V,A)\geq1\), a contradiction.

Thus \(M=0\).  The section \(s_0\) is nowhere zero.  Indeed, if a local
representative were a nonunit at a point of the normal space, a minimal
prime over its nonzero principal ideal would have height one, and would
give a component of its zero divisor.  Consequently \(s_0\) supplies
an isomorphism of holomorphic lines
\[
 \OO_X\xrightarrow{\ \simeq\ }\OO_X(m_0D).
\]
This proves the asserted actual \(\Q\)-linear triviality in Iitaka
dimension zero, and in particular semiampleness there.  Together with
the positive-Iitaka-dimensional case it proves the theorem.
\end{proof}
